\chapter{Removing the Wastes}\label{ch:g7:removing-wastes}

Deliveries have dominated this year: \omterm{def:g7:what-is-respiration:respiration}{oxygen} in, \omterm{def:g7:digestion-and-nutrients:families}{nutrients}
in, fuel to every working \omterm{def:g5:first-look-at-cells:cell}{cell}. But every workshop makes
waste, and a delivery service that never collects the
rubbish soon poisons its own city. The \omterm{def:g7:what-is-respiration:respiration}{carbon dioxide}
account is already settled --- breathed out at the \omterm{def:g7:lungs-and-gas-exchange:alveoli}{alveoli}.
This closing chapter follows the \emph{other} wastes, to
the body's discreetest organs: the \omterm{def:g7:removing-wastes:kidneys}{kidneys}.

\section{The wastes beyond carbon dioxide}

\begin{proposition}[The cells' other rubbish]\label{prop:g7:removing-wastes:rubbish}
The \omterm{def:g5:first-look-at-cells:cell}{cells}' burning and building
(\cref{prop:g7:what-is-respiration:power},
\cref{prop:g6:matter-from-food:transform}) yield more than
\omterm{def:g7:what-is-respiration:respiration}{carbon dioxide}. Using building-matter \omterm{def:g7:digestion-and-nutrients:families}{nutrients}, in
particular, \omterm{def:g1:plants-around-us:parts}{leaves} a nitrogen-bearing waste ---
\emph{urea}\index{urea} --- which the \omterm{prop:g4:journey-of-food:absorb}{blood} collects from
every tissue exactly as it collects \omterm{def:g7:what-is-respiration:respiration}{carbon dioxide}. \omterm{prop:g7:removing-wastes:rubbish}{Urea}
is harmless in passing traffic and poisonous piled up: it
must leave, continuously.
\end{proposition}
\admitted

\begin{example}[Two collection rounds, one blood]\label{ex:g7:removing-wastes:rounds}
The same drop of \cref{ex:g7:heart-and-circulation:tour}
runs both collections at once: at the thigh's \omterm{def:g7:heart-and-circulation:vessels}{capillaries}
it takes on \omterm{def:g7:what-is-respiration:respiration}{carbon dioxide} \emph{and} \omterm{prop:g7:removing-wastes:rubbish}{urea}. The first it
drops at the \omterm{def:g7:how-animals-breathe:four}{lungs}, next circuit. The second rides on ---
past the \omterm{def:g7:how-animals-breathe:four}{lungs}, which cannot touch it --- until the body
circuit carries it through the organs built for exactly
this cargo.
\end{example}

\section{The kidneys, the blood's filters}

\begin{definition}[Kidneys and the urinary system]\label{def:g7:removing-wastes:kidneys}
The two \emph{kidneys}\index{kidney} --- fist-sized,
bean-shaped, low in the back --- are blood-cleaning
filters. Each receives a generous \omterm{def:g7:heart-and-circulation:vessels}{artery} straight from
the body circuit's main line. Inside, the \omterm{prop:g4:journey-of-food:absorb}{blood} is
filtered through about a million microscopic filter
units; the cleaned \omterm{prop:g4:journey-of-food:absorb}{blood} returns by the kidney's \omterm{def:g7:heart-and-circulation:vessels}{vein},
and the removed load --- \omterm{prop:g7:removing-wastes:rubbish}{urea}, surplus water, surplus
salts --- \omterm{def:g1:plants-around-us:parts}{leaves} as \emph{urine}\index{urine}, trickling
down two tubes to the \emph{bladder}\index{bladder}, the
holding tank emptied at will.
\end{definition}

\begin{proposition}[Filter, then reclaim]\label{prop:g7:removing-wastes:filter}
The \omterm{def:g7:removing-wastes:kidneys}{kidney}'s trick is generosity followed by thrift. The
filter units first push a huge volume of the \omterm{prop:g4:journey-of-food:absorb}{blood}'s
watery part --- many tens of litres a day --- through
their fine sieves, \omterm{prop:g7:removing-wastes:rubbish}{urea}, sugars, salts and all. Then,
along each unit's winding tube, the body \emph{reclaims}
what it wants back: nearly all the water, all the
\omterm{prop:g7:organs-at-work:bill}{glucose}, salts to measure --- leaving the \omterm{prop:g7:removing-wastes:rubbish}{urea} and the
surpluses behind. Of the tens of litres sieved, about a
litre and a half \omterm{def:g1:plants-around-us:parts}{leaves} as \omterm{def:g7:removing-wastes:kidneys}{urine}: filter everything,
take back the good, and the waste has nowhere to hide.
\end{proposition}
\admitted

\begin{omfigure}
\begin{tikzpicture}[scale=0.9]
  % blood in
  \draw[->, very thick, omThm] (-0.2,3.0) -- (2.2,3.0);
  \node[font=\footnotesize, omThm, above, align=center] at (1.0,3.1)
    {blood in\\ (with urea)};
  % kidney box
  \draw[thick, fill=omDef!6, rounded corners=8pt]
    (2.3,1.0) rectangle (8.1,4.6);
  \node[font=\small\bfseries, omDef] at (5.2,4.25) {the kidney};
  % filter stage
  \draw[thick, omDef, fill=white] (2.7,2.2) rectangle (4.5,3.6);
  \node[font=\footnotesize, align=center] at (3.6,2.9)
    {sieve:\\ filter all};
  % reclaim stage
  \draw[thick, omMeth, fill=white] (5.0,2.2) rectangle (7.7,3.6);
  \node[font=\footnotesize, align=center] at (6.35,2.9)
    {reclaim:\\ water, glucose,\\ salts to measure};
  \draw[->, thick] (4.5,2.9) -- (5.0,2.9);
  % blood out
  \draw[->, very thick, omDef] (8.1,3.0) -- (10.5,3.0);
  \node[font=\footnotesize, omDef, above, align=center] at (9.3,3.1)
    {cleaned\\ blood out};
  % urine out
  \draw[->, very thick, omProp] (6.35,2.2) -- (6.35,0.0);
  \node[font=\footnotesize, omProp, right, align=left] at (6.5,0.55)
    {urine: urea, surplus\\ water and salts};
  \node[font=\footnotesize, align=center] at (6.35,-0.35)
    {to the bladder};
\end{tikzpicture}

{\small The \omterm{def:g7:removing-wastes:kidneys}{kidney}'s two-stage trick: sieve generously,
reclaim thriftily. What is not taken back --- \omterm{prop:g7:removing-wastes:rubbish}{urea} and the
surpluses --- is the \omterm{def:g7:removing-wastes:kidneys}{urine}.}
\end{omfigure}

\begin{example}[Reading the thrift]\label{ex:g7:removing-wastes:thrift}
The reclaiming stage explains the everyday observations.
Drink much, and the \omterm{def:g7:removing-wastes:kidneys}{kidneys} reclaim less water: \omterm{def:g7:removing-wastes:kidneys}{urine}
plentiful and pale. \omterm{rem:g7:removing-wastes:sweat}{Sweat} a hot match dry, and they
reclaim nearly all: \omterm{def:g7:removing-wastes:kidneys}{urine} scant and dark --- the body
guarding its water
(\cref{prop:g4:what-plants-need:needs} taught water's
value for \omterm{def:g1:plants-around-us:plant}{plants}; the \omterm{def:g7:removing-wastes:kidneys}{kidney} is the \omterm{def:g1:animals-around-us:animal}{animal} side of the
same economy). And \omterm{prop:g7:organs-at-work:bill}{glucose} in the \omterm{def:g7:removing-wastes:kidneys}{urine}, which the sieve
lets through but reclaiming should recover completely, is
a classic medical signal that something --- reclaiming or
\omterm{prop:g4:journey-of-food:absorb}{blood} sugar --- is out of order.
\end{example}

\begin{remark}[The skin lends a hand]\label{rem:g7:removing-wastes:sweat}
\emph{Sweat}\index{sweat} is a sideline of the same
business: water and salts (with traces of \omterm{prop:g7:removing-wastes:rubbish}{urea}) released
at the \omterm{def:g1:the-five-senses:sense}{skin} --- though its first job is
\cref{ex:g7:organs-at-work:felt}'s cooling. The \omterm{def:g7:how-animals-breathe:four}{lungs},
\omterm{def:g7:removing-wastes:kidneys}{kidneys} and \omterm{def:g1:the-five-senses:sense}{skin} divide the body's whole waste removal:
\omterm{def:g7:what-is-respiration:respiration}{carbon dioxide} to the first, \omterm{prop:g7:removing-wastes:rubbish}{urea} and the surpluses to
the second, heat and a salty trickle to the third.
Nothing poisonous is left to pile up --- in a healthy
body, the rubbish never wins.
\end{remark}

\section{The year's machine, complete}

\begin{proposition}[The full supply-and-removal loop]\label{prop:g7:removing-wastes:complete}
With the \omterm{def:g7:removing-wastes:kidneys}{kidneys} in place, this year's tour closes into
one machine. The \omterm{prop:g4:journey-of-food:absorb}{blood}, pumped round the double circuit
(\cref{def:g7:heart-and-circulation:double}), loads
\omterm{def:g7:what-is-respiration:respiration}{oxygen} at the \omterm{def:g7:lungs-and-gas-exchange:alveoli}{alveoli} and \omterm{def:g7:digestion-and-nutrients:families}{nutrients} at the \omterm{prop:g7:digestion-and-nutrients:absorption}{villi};
delivers both to every working \omterm{def:g5:first-look-at-cells:cell}{cell} for burning and
building; and collects the wastes --- \omterm{def:g7:what-is-respiration:respiration}{carbon dioxide} to
the \omterm{def:g7:how-animals-breathe:four}{lungs}, \omterm{prop:g7:removing-wastes:rubbish}{urea} and surpluses to the \omterm{def:g7:removing-wastes:kidneys}{kidneys}. Every organ
of grades four and seven is a station of this single
loop; every exchange crosses a \omterm{def:g7:heart-and-circulation:vessels}{capillary} wall; and the
whole runs on \cref{prop:g7:what-is-respiration:power}'s
slow fire, \omterm{def:g5:first-look-at-cells:cell}{cell} by \omterm{def:g5:first-look-at-cells:cell}{cell}.
\end{proposition}
\admitted

\begin{method}[The whole-body audit]\label{met:g7:removing-wastes:audit}
For any substance, audit its passage:
\begin{enumerate}
  \item entry: mouth and \omterm{prop:g7:digestion-and-nutrients:absorption}{villi}, or nose and \omterm{def:g7:lungs-and-gas-exchange:alveoli}{alveoli};
  \item transport: which circuit \omterm{def:g1:my-body:parts}{legs}, which special
        stops (\omterm{prop:g7:digestion-and-nutrients:absorption}{liver});
  \item use: fuel burned, material built in, or neither;
  \item exit: \omterm{def:g7:lungs-and-gas-exchange:alveoli}{alveoli} (gas), \omterm{def:g7:removing-wastes:kidneys}{kidneys} (\omterm{prop:g7:removing-wastes:rubbish}{urea}, surpluses),
        \omterm{def:g1:the-five-senses:sense}{skin} (sweat's share) --- or, never absorbed at
        all, the leftovers' route
        (\cref{prop:g4:journey-of-food:absorb}).
\end{enumerate}
Run it on water, on \omterm{prop:g7:organs-at-work:bill}{glucose}, on a salty crisp: the
machine has a station for every step, which is the
year's finding in one exercise.
\end{method}

\section{Exercises}

\begin{exercise}[$\star$]\label{exo:g7:removing-wastes:1}
What is \omterm{prop:g7:removing-wastes:rubbish}{urea} --- from what use does it arise, and why
must it leave?
\end{exercise}

\begin{exercise}[$\star$]\label{exo:g7:removing-wastes:2}
Which wastes leave at the \omterm{def:g7:how-animals-breathe:four}{lungs}, which at the \omterm{def:g7:removing-wastes:kidneys}{kidneys},
and which at the \omterm{def:g1:the-five-senses:sense}{skin}?
\end{exercise}

\begin{exercise}[$\star$]\label{exo:g7:removing-wastes:3}
Describe the \omterm{def:g7:removing-wastes:kidneys}{kidneys}: number, size, place, and their
plumbing to the \omterm{def:g7:removing-wastes:kidneys}{bladder}.
\end{exercise}

\begin{exercise}[$\star$]\label{exo:g7:removing-wastes:4}
State the filter-then-reclaim trick, with the two
volumes that frame it.
\end{exercise}

\begin{exercise}[$\star$]\label{exo:g7:removing-wastes:5}
What is reclaimed completely, what to measure, and what
not at all?
\end{exercise}

\begin{exercise}[$\star$]\label{exo:g7:removing-wastes:6}
Why is \omterm{def:g7:removing-wastes:kidneys}{urine} plentiful and pale one day, scant and dark
another?
\end{exercise}

\begin{exercise}[$\star\star$]\label{exo:g7:removing-wastes:7}
Why can the \omterm{def:g7:how-animals-breathe:four}{lungs} not remove \omterm{prop:g7:removing-wastes:rubbish}{urea}, and the \omterm{def:g7:removing-wastes:kidneys}{kidneys} not
remove \omterm{def:g7:what-is-respiration:respiration}{carbon dioxide} fast enough? What does the divided
labor say about the two wastes' forms?
\end{exercise}

\begin{exercise}[$\star\star$]\label{exo:g7:removing-wastes:8}
Explain why \omterm{prop:g7:organs-at-work:bill}{glucose} in the \omterm{def:g7:removing-wastes:kidneys}{urine} alarms a doctor,
through the sieve-and-reclaim stages.
\end{exercise}

\begin{exercise}[$\star\star$]\label{exo:g7:removing-wastes:9}
Follow one \omterm{prop:g7:removing-wastes:rubbish}{urea} unit from a thigh \omterm{def:g3:bones-and-muscles:muscle}{muscle}'s \omterm{def:g7:heart-and-circulation:vessels}{capillary} to
the \omterm{def:g7:removing-wastes:kidneys}{bladder}, station by station.
\end{exercise}

\begin{exercise}[$\star\star$]\label{exo:g7:removing-wastes:10}
Run \cref{met:g7:removing-wastes:audit} on a glass of
water drunk after sport.
\end{exercise}

\begin{exercise}[$\star\star$]\label{exo:g7:removing-wastes:11}
\omterm{def:g7:removing-wastes:kidneys}{Kidney} patients whose filters fail visit a machine that
cleans their \omterm{prop:g4:journey-of-food:absorb}{blood} through an artificial \omterm{def:g6:cell-unit-of-life:parts}{membrane}.
Which stage of \cref{prop:g7:removing-wastes:filter}
does the machine copy, and why must sessions recur
every few days?
\end{exercise}

\begin{exercise}[$\star\star\star$]\label{exo:g7:removing-wastes:12}
Write the year's summary paragraph promised by
\cref{prop:g7:removing-wastes:complete}: one \omterm{prop:g4:journey-of-food:absorb}{blood}, two
circuits, four borders (\omterm{def:g7:lungs-and-gas-exchange:alveoli}{alveoli}, \omterm{prop:g7:digestion-and-nutrients:absorption}{villi}, \omterm{def:g3:bones-and-muscles:muscle}{muscle}
\omterm{def:g7:heart-and-circulation:vessels}{capillaries}, \omterm{def:g7:removing-wastes:kidneys}{kidney} filters) --- and the slow fire that
the whole machine serves.
\end{exercise}

\section{Problem: The Day of a Water Molecule}

\begin{problem}[{Weekend problem --- one glass of water,
audited through the whole machine}]\label{pb:g7:removing-wastes:1}
Follow a glass of water drunk at breakfast on a sports
day, as the year's grand audit.

\medskip\noindent\textbf{Part I --- Entry and transport.}
\begin{enumerate}
  \item Water needs no dismantling
        (\cref{def:g7:digestion-and-nutrients:families}).
        Name its entry crossing and the vessel network
        that receives it.
  \item Its \omterm{prop:g4:journey-of-food:absorb}{blood} makes the standard first stop. Which
        organ, and why is the routing sensible for
        \omterm{def:g7:digestion-and-nutrients:families}{nutrients} even if water itself needs no sorting?
  \item By mid-morning the water is everywhere the \omterm{prop:g4:journey-of-food:absorb}{blood}
        goes. Name the two circuits it now rides in
        alternation.
  \item Some of it is soon inside working \omterm{def:g5:first-look-at-cells:cell}{cells}. List
        two jobs water does there (recall that the
        \omterm{def:g5:first-look-at-cells:cell}{cells}' chemistry --- juices, burning --- runs
        wet).
\end{enumerate}

\medskip\noindent\textbf{Part II --- The afternoon
match.}
\begin{enumerate}[resume]
  \item During the match, part of the water \omterm{def:g1:plants-around-us:parts}{leaves} by a
        border that is not an excretory organ first.
        Name it, its cargo besides water, and its real
        job that day.
  \item The match over, the player is thirsty and her
        \omterm{def:g7:removing-wastes:kidneys}{urine} that evening is scant and dark. Explain
        both with the reclaiming stage.
  \item Meanwhile the match's hard-worked \omterm{def:g3:bones-and-muscles:muscle}{muscles} made
        extra \omterm{def:g7:what-is-respiration:respiration}{carbon dioxide} and \omterm{prop:g7:removing-wastes:rubbish}{urea}. Assign each its
        exit border, and the circuit \omterm{def:g1:my-body:parts}{legs} that took it
        there.
  \item Why did the \omterm{def:g7:removing-wastes:kidneys}{kidneys} keep working right through
        the match, though \omterm{def:g7:removing-wastes:kidneys}{urine} production fell? What
        were they still removing?
\end{enumerate}

\medskip\noindent\textbf{Part III --- The audit closed.}
\begin{enumerate}[resume]
  \item Next morning, some of the glass has left in
        breath. Explain how water exits at the \omterm{def:g7:lungs-and-gas-exchange:alveoli}{alveoli}
        (\cref{ex:g4:breathing:mirror} showed it years
        ago).
  \item Draw up the glass's full ledger: entry, stores
        and uses, and its four exits --- \omterm{def:g7:removing-wastes:kidneys}{urine}, sweat,
        breath, and the small share in the leftovers.
  \item The audit shows one \omterm{prop:g4:journey-of-food:absorb}{blood} serving \omterm{def:g7:how-animals-breathe:four}{lungs}, \omterm{prop:g7:digestion-and-nutrients:absorption}{villi},
        \omterm{def:g3:bones-and-muscles:muscle}{muscles}, \omterm{def:g7:removing-wastes:kidneys}{kidneys} and \omterm{def:g1:the-five-senses:sense}{skin} in turn. State, in
        two sentences, why this single-network design
        makes every organ's health every other organ's
        business.
  \item Close the year: which single idea --- name it
        --- did every station of this audit turn out to
        serve? (\cref{prop:g7:what-is-respiration:power}
        holds the answer.)
\end{enumerate}
\end{problem}
